\chapter{Υλοποίηση}
\label{ch:implementation}

\section{Εργαλεία και τεχνολογίες}
\label{sec:impl-tools}

Η προσομοίωση link-budget υλοποιείται εξ ολοκλήρου σε MATLAB, αξιοποιώντας
τρία πακέτα εργαλείων (toolboxes):

\begin{itemize}
  \item Το \textbf{5G Toolbox}~\cite{matlab5gtoolbox} (\texttt{nrCarrierConfig},
    \texttt{nrPathLossConfig}, \texttt{nrPathLoss}), για τη διαμόρφωση του
    φέροντος NR και τον υπολογισμό path loss κατά TR~38.901.
  \item Το \textbf{Mapping Toolbox} (\texttt{wgs84Ellipsoid},
    \texttt{geodetic2enu}, \texttt{geodetic2aer}, \texttt{distance}), για
    γεωδαιτικούς υπολογισμούς (μετατροπές συντεταγμένων, αποστάσεις, γωνία
    ανύψωσης).
  \item Το \textbf{Communications/Phased Array Toolbox}
    (\texttt{fspl}, \texttt{physconst}), για τον υπολογισμό απωλειών
    ελεύθερου χώρου και φυσικών σταθερών (ταχύτητα φωτός, σταθερά Boltzmann).
\end{itemize}

Το δεύτερο, μεταγενέστερο σκέλος -- η δοκιμή μηχανικής μάθησης -- υλοποιείται
σε Python, με τις βιβλιοθήκες pandas (διαχείριση δεδομένων),
scikit-learn~\cite{pedregosa2011sklearn} (ταξινομητές, προεπεξεργασία,
μετρικές αξιολόγησης), και matplotlib (γραφήματα).

\section{Δομή λογισμικού}
\label{sec:impl-structure}

Ο κώδικας οργανώνεται σε δύο κύριους φακέλους:

\begin{description}
  \item[\texttt{PROD/}] Ο πυρήνας της προσομοίωσης και οι driver scripts.
  \begin{description}
    \item[\texttt{simulateScenario.m}] Ο πυρήνας υπολογισμού: δέχεται
      γεωμετρία (θέσεις BS/χρηστών/δορυφόρου) και ραδιο-παραμέτρους,
      επιστρέφει per-user KPI και per-candidate διαγνωστικά. Καλείται από
      όλα τα υπόλοιπα scripts -- είναι η μοναδική υλοποίηση του link-budget,
      ώστε στατικό σενάριο, Monte Carlo dataset, στατιστική ανάλυση, και
      χρονική προσομοίωση να χρησιμοποιούν \emph{ακριβώς} το ίδιο μοντέλο
      καναλιού.
    \item[\texttt{test\_simulation.m}] Το σενάριο αναφοράς: ορίζει τη
      στατική τοπολογία και τις ραδιο-παραμέτρους, καλεί μία φορά τη
      \texttt{simulateScenario}, και εμφανίζει πίνακα αποτελεσμάτων
      (\texttt{array.m}) και τρισδιάστατη οπτικοποίηση (\texttt{visual.m}).
    \item[\texttt{monteCarloDriver.m}] Παράγει το ετικετοποιημένο σύνολο
      δεδομένων για το Μέρος~2 (μηχανική μάθηση), καλώντας τη
      \texttt{simulateScenario} πάνω σε τυχαιοποιημένες τοπολογίες
      (Ενότητα~\ref{sec:meth-montecarlo}).
    \item[\texttt{kpiRepeatedRuns.m}] Επαναλαμβάνει το \emph{ίδιο} στατικό
      σενάριο με διαφορετικό RNG seed ανά εκτέλεση, ώστε να ποσοτικοποιηθεί
      η διακύμανση KPI λόγω στοχαστικότητας καναλιού
      (Ενότητα~\ref{sec:results-repeated}).
    \item[\texttt{temporalPassSimulation.m}] Η χρονική προσομοίωση πάσσου
      LEO (Ενότητα~\ref{sec:meth-temporal}, αποτελέσματα στην
      Ενότητα~\ref{sec:results-temporal}).
  \end{description}
  \item[\texttt{Model/}] Το δεύτερο σκέλος (Python).
  \begin{description}
    \item[\texttt{train\_model.py}] Φορτώνει το CSV που παράγει ο
      \texttt{monteCarloDriver.m}, εκπαιδεύει τους δύο ταξινομητές
      (Ενότητα~\ref{sec:meth-ml}), και αποθηκεύει μετρικές/γραφήματα.
  \end{description}
\end{description}

\section{Ροή δεδομένων}
\label{sec:impl-dataflow}

Η συνολική ροή δεδομένων του pipeline απεικονίζεται σχηματικά ως εξής:

\begin{center}
\small
\texttt{γεωγραφικές συντεταγμένες + ραδιο-παράμετροι}
$\rightarrow$ \texttt{simulateScenario.m}
$\rightarrow$ \texttt{per-user KPI (MATLAB table)}
$\rightarrow$ \texttt{CSV (writetable)}
$\rightarrow$ \texttt{pandas DataFrame}
$\rightarrow$ \texttt{scikit-learn pipeline}
$\rightarrow$ \texttt{μετρικές (JSON) + γραφήματα (PNG)}
\end{center}

Η ενδιάμεση αποθήκευση σε CSV -- αντί, π.χ., για απευθείας κλήση Python
μέσα από τη MATLAB -- κρατά τα δύο σκέλη (προσομοίωση link-budget σε
MATLAB, εκπαίδευση μοντέλου σε Python) πλήρως ανεξάρτητα και εύκολα
ελέγξιμα το καθένα ξεχωριστά, με το CSV να λειτουργεί ως σαφές,
ανθρωπίνως αναγνώσιμο συμβόλαιο μεταξύ τους.

\section{Αναπαραγωγιμότητα}
\label{sec:impl-reproducibility}

Όλα τα scripts καλούν \texttt{rng(seed)} με ρητή τιμή σπόρου πριν από κάθε
στοχαστική εκτέλεση: σταθερό seed (\texttt{rng(42)}) για το στατικό
σενάριο αναφοράς και τον Monte Carlo dataset generator, και διαφορετικό,
αλλά ντετερμινιστικά καθορισμένο seed ανά επανάληψη (\texttt{rng(r)} για
το $r$-οστό run, \texttt{rng(step+1)} για το $\text{step}$-οστό χρονικό
βήμα) στα \texttt{kpiRepeatedRuns.m} και \texttt{temporalPassSimulation.m}
αντίστοιχα. Αυτό εξασφαλίζει ότι κάθε αριθμητικό αποτέλεσμα που
παρουσιάζεται στο Κεφάλαιο~\ref{ch:results} είναι πλήρως αναπαραγώγιμο.
